\documentclass[12pt,a4paper]{article}

% --- Настройка языка и шрифтов для LuaLaTeX ---
\usepackage{fontspec}
\defaultfontfeatures{Ligatures=TeX,Scale=MatchLowercase}

\setmainfont{Alice-Regular}[
    Path=font/,
    Extension=.ttf,
    BoldFont = Alice-Regular,
    ItalicFont = Alice-Regular,
    BoldItalicFont = Alice-Regular,
    BoldFeatures = {FakeBold=1.6},
    ItalicFeatures = {FakeSlant=0.2},
    BoldItalicFeatures = {FakeBold=1.6,FakeSlant=0.2}
]

\newfontfamily\cyrillicfont{Alice-Regular}[
    Path=font/,
    Extension=.ttf,
    BoldFont = Alice-Regular,
    ItalicFont = Alice-Regular,
    BoldItalicFont = Alice-Regular,
    BoldFeatures = {FakeBold=1.6},
    ItalicFeatures = {FakeSlant=0.2},
    BoldItalicFeatures = {FakeBold=1.6,FakeSlant=0.2}
]

\usepackage{polyglossia}
\setmainlanguage{russian}
\setsansfont{TeX Gyre Heros}
\setmonofont{TeX Gyre Cursor}

\usepackage{graphicx}        					% Графика
\usepackage{amsmath, amsfonts, amssymb, amsthm}	% Математика
\usepackage{mathtools}
\usepackage{chemfig}
\usepackage[version=4,arrows=pgf-filled,
            textfontname=sffamily,
            mathfontname=mathsf]{mhchem}
\usepackage[protrusion=false]{microtype}       % Микротипографика
\usepackage{braket} 						% Braket нотация для квантовой механики
\usepackage{unicode-math} 					% Современный пакет для математики в LuaLaTeX
\usepackage{longtable}       				% Таблицы с разрывом страниц
\usepackage{tikz}           				% Рисование
\usepackage[europeanresistors]{circuitikz}	% Рисование электрических схем в европейском стиле
\usepackage{listings}						% Для отображения кода
\usepackage{tabularx}       				% Для таблиц во всю ширину
\usepackage{array}		  					% Для расширенных возможностей таблиц
\usepackage{cancel}                % Перечёркивать
\usepackage{float}          				% Для точного позиционирования фигур
\usepackage{pgfplots}                          % Для построения графиков
\usepackage{caption}                           % Для настройки подписей к рисункам и таблицам
\usepackage{xfp}                               % Для вычислений в LaTeX
\usepackage{luacode}                           % Для выполнения Lua-кода в документе    
\pgfplotsset{compat=1.18}
\usepackage[paper=letterpaper,
            top=20mm,
            bottom=20mm,
            left=10mm,
            right=10mm,
            footskip=0.4in 
]{geometry}
\usepackage{colortbl}

\usetikzlibrary{arrows.meta}

% --- Колонтитулы ---
\usepackage{fancyhdr}
\pagestyle{fancy}
\fancyhf{}

% --- Теоремы и окружения ---
\theoremstyle{plain}
\newtheorem{theorem}{Theorem}[section]
\newtheorem{lemma}[theorem]{Lemma}
\newtheorem{proposition}[theorem]{Proposition}
\newtheorem{corollary}[theorem]{Corollary}

\theoremstyle{definition}
\newtheorem{definition}[theorem]{Definition}

\theoremstyle{remark}
\newtheorem{remark}[theorem]{Remark}

\numberwithin{equation}{section}


\renewcommand{\headrulewidth}{0.4pt}
\renewcommand{\footrulewidth}{0.4pt}
\renewcommand\lstlistingname{Исходный код}

\lstdefinestyle{main}{
	backgroundcolor=\color{gray!8},   
	commentstyle=\color{green!50!black},
	keywordstyle=\color{blue!80!black},
	numberstyle=\small\color{darkgray},
	stringstyle=\rmfamily\color{orange!90!black},
	basicstyle=\ttfamily\small,
	breakatwhitespace=false,         
	breaklines=true,                   
	keepspaces=true,                 
	numbers=left,                    
	numbersep=8pt,                  
	showspaces=false,                
	showstringspaces=false,
	showtabs=false,                  
	tabsize=4
}	

\lstset{style=main}

\newcolumntype{C}{>{\centering\arraybackslash}X}
\renewcommand{\arraystretch}{1.25}

% =========================================================
%         СЛУЖЕБНЫЕ КОМАНДЫ ДЛЯ ВЫЧИСЛЕНИЙ
% =========================================================
\newcommand{\CalcZero} [1]{\fpeval{round((#1),0)}}
\newcommand{\CalcOne}  [1]{\fpeval{round((#1),1)}}
\newcommand{\CalcTwo}  [1]{\fpeval{round((#1),2)}}
\newcommand{\CalcThree}[1]{\fpeval{round((#1),3)}}
\newcommand{\CalcFour} [1]{\fpeval{round((#1),4)}}
\newcommand{\CalcFive} [1]{\fpeval{round((#1),5)}}

\begin{luacode*}
-- безопасная версия value_with_mantice
local function round_fallback(x)
  if x >= 0 then return math.floor(x + 0.5) else return math.ceil(x - 0.5) end
end
local round = math.round or round_fallback

function value_with_mantice(value, precision, shift, scale)
  precision = precision or 5
  shift = shift or 0
  scale = scale or 0

  if value == nil or value ~= value then
    return "\\fpeval{0}"
  end
  if value == 0 then
    return "\\fpeval{0}"
  end

  local sign = ""
  if value < 0 then
    sign = "-"
    value = math.abs(value)
  end

  local logv = math.log10(value)
  if not (logv == logv) then -- проверка на NaN
    return "\\fpeval{0}"
  end

  local floorlog = math.floor(logv)
  local mantissa = round(value * 10^(precision - 1 - floorlog)) * 10^(-(precision - 1) + shift - scale)
  local exponent = floorlog - shift

  if exponent == 0 then
    return sign .. "\\fpeval{" .. mantissa .. "}"
  else
    return sign .. "\\fpeval{" .. mantissa .. "} \\cdot 10^{" .. exponent .. "}"
  end
end

\end{luacode*}

\begin{document}

\section{Критерий Поста}

\noindent Полный набор функций от двух переменных

\vspace{0.5em}

\noindent\begin{tabularx}{\linewidth}{|C|C|C|C|C|C|C|C|C|C|C|c|C|c|C|c|C|C|}
  \hline
  $x$ & $y$ & $\phi_0$ & $\land$ & $\phi_2$ & $\phi_3$ & $\phi_4$ & $\phi_5$ & $+$ & $\lor$ & $\downarrow$ & $\iff$ & $\phi_{10}$ & $\impliedby$ & $\phi_{12}$ & $\implies$ & $|$ & $\phi_{15}$ \\
  \hline
  $0$ & $0$ & $0$ & $0$ & $0$ & $0$ & $0$ & $0$ & $0$ & $0$ & $1$ & $1$ & $1$ & $1$ & $1$ & $1$ & $1$ & $1$ \\
  $0$ & $1$ & $0$ & $0$ & $0$ & $0$ & $1$ & $1$ & $1$ & $1$ & $0$ & $0$ & $0$ & $0$ & $1$ & $1$ & $1$ & $1$ \\
  $1$ & $0$ & $0$ & $0$ & $1$ & $1$ & $0$ & $0$ & $1$ & $1$ & $0$ & $0$ & $1$ & $1$ & $0$ & $0$ & $1$ & $1$ \\
  $1$ & $1$ & $0$ & $1$ & $0$ & $1$ & $0$ & $1$ & $0$ & $1$ & $0$ & $1$ & $0$ & $1$ & $0$ & $1$ & $0$ & $1$ \\
  \hline
\end{tabularx} 

\vspace{0.5em}

\noindent Классы Поста для функций двух переменных

\vspace{0.5em}

\noindent\begin{tabularx}{\linewidth}{|c|c|c|c|c|c|X|}
  \hline
               & $T_0$ & $T_1$ & $S$ & $M$ & $L$ & Мн. Жегалкина \\
  \hline
  $\phi_0$     &  $1$  &  $0$  & $0$ & $1$ & $1$ & $0$ \\
  $\land$      &  $1$  &  $1$  & $0$ & $1$ & $0$ & $xy$ \\
  $\phi_2$     &  $1$  &  $0$  & $0$ & $0$ & $0$ & $x + xy$ \\
  $\phi_3$     &  $1$  &  $1$  & $1$ & $1$ & $1$ & $x$ \\
  $\phi_4$     &  $1$  &  $0$  & $0$ & $0$ & $0$ & $y + xy$ \\
  $\phi_5$     &  $1$  &  $1$  & $1$ & $1$ & $1$ & $y$ \\
  $+$          &  $1$  &  $0$  & $0$ & $0$ & $1$ & $x + y$ \\
  $\lor$       &  $1$  &  $1$  & $0$ & $1$ & $0$ & $x + y + xy$ \\
  $\downarrow$ &  $0$  &  $0$  & $0$ & $0$ & $0$ & $1 + x + y + xy$ \\
  $\iff$       &  $0$  &  $1$  & $0$ & $0$ & $1$ & $1 + x + y$ \\
  $\phi_{10}$  &  $0$  &  $0$  & $1$ & $0$ & $1$ & $1 + y$ \\
  $\impliedby$ &  $0$  &  $1$  & $0$ & $0$ & $0$ & $1 + y + xy$ \\
  $\phi_{12}$  &  $0$  &  $0$  & $1$ & $0$ & $1$ & $1 + x$ \\
  $\implies$   &  $0$  &  $1$  & $0$ & $0$ & $0$ & $1 + x + xy$ \\
  $\phi_{14}$  &  $0$  &  $0$  & $0$ & $0$ & $0$ & $1 + xy$ \\
  $\phi_{15}$  &  $0$  &  $1$  & $0$ & $1$ & $1$ & $1$ \\
  \hline
\end{tabularx}

\subsection{Образуют ли системы булевых функций базис?}
\begin{enumerate}
  \renewcommand{\labelenumi}{\asbuk{enumi})}
  \item $\mathcal{B} = \{\phi_0,\,\phi_1,\,\phi_{15}\}$
  \item $\mathcal{B} = \{\phi_2,\,\phi_{10},\,\phi_{11}\}$
  \item $\mathcal{B} = \{\phi_6,\,\phi_7,\,\phi_{9}\}$
\end{enumerate}

\subsubsection{Образует ли система булевых функций а) базис?}

\noindent\begin{tabularx}{\linewidth}{|c|c|c|c|c|c|X|}
  \hline
               & $T_0$ & $T_1$ & \textcolor{red}{$S$} & \textcolor{red}{$M$} & $L$ & Мн. Жегалкина \\
  \hline
  $\phi_0$     &  $1$  &  $0$  & $0$ & $1$ & $1$ & $0$ \\
  $\land$      &  $1$  &  $1$  & $0$ & $1$ & $0$ & $xy$ \\
  $\phi_{15}$  &  $0$  &  $1$  & $0$ & $1$ & $1$ & $1$ \\
  \hline
\end{tabularx}

\vspace{0.5em}
\noindent Ответ: система не образует базис.

\subsubsection{Образует ли система булевых функций б) базис?}

\noindent\begin{tabularx}{\linewidth}{|c|c|c|c|c|c|X|}
  \hline
               & \textcolor{red}{$T_0$} & \textcolor{red}{$T_1$} & \textcolor{red}{$S$} & \textcolor{red}{$M$} & \textcolor{red}{$L$} & Мн. Жегалкина \\
  \hline
  $\phi_2$     &  $1$  &  $0$  & $0$ & $0$ & $0$ & $x + xy$ \\
  $\phi_{10}$  &  $0$  &  $0$  & $1$ & $0$ & $1$ & $1 + y$ \\
  $\impliedby$ &  $0$  &  $1$  & $0$ & $0$ & $0$ & $1 + y + xy$ \\
  \hline
\end{tabularx}

\vspace{0.5em}
\noindent Ответ: система не образует базис.

\subsubsection{Образует ли система булевых функций в) базис?}

\noindent\begin{tabularx}{\linewidth}{|c|c|c|c|c|c|X|}
  \hline
               & $T_0$ & $T_1$ & \textcolor{red}{$S$} & \textcolor{red}{$M$} & $L$ & Мн. Жегалкина \\
  \hline
  $+$          &  $1$  &  $0$  & $0$ & $0$ & $1$ & $x + y$ \\
  $\lor$       &  $1$  &  $1$  & $0$ & $1$ & $0$ & $x + y + xy$ \\
  $\iff$       &  $0$  &  $1$  & $0$ & $0$ & $1$ & $1 + x + y$ \\
  \hline
\end{tabularx}

\vspace{0.5em}
\noindent Ответ: система не образует базис.

\subsection{Найдите все базисы, состоящие из двух булевых функций двух переменных.}

Так как необходимо из двух булевых функций получить по одному исключению каждого 
класса Поста, то исключаются функции самобазисы $\downarrow$ и $\phi_{14}$. Функции
$\phi_3$ и $\phi_5$ принадлежат всем классам Поста, поэтому с ними невозможно
образовать минимальную систему, какой является базис. Для функций $\phi_0$, 
$\phi_2$, $\phi_4$, $+$, $\iff$, $\phi_{15}$, $\implies$ и $\impliedby$ нет пар 
функций. Остаются $\land$ и $\lor$ с $\phi_{10}$ и $\phi_{12}$.  
\begin{enumerate}
  \item $\mathcal{B} = \{\phi_{10},\,\land\}$
  \item $\mathcal{B} = \{\phi_{10},\,\lor\}$
  \item $\mathcal{B} = \{\phi_{12},\,\land\}$
  \item $\mathcal{B} = \{\phi_{12},\,\lor\}$
\end{enumerate}  

\subsection{Найдите все базисы, состоящие их трёх булевых функций двух переменных.}

Так как необходимо из двух булевых функций получить по одному исключению каждого 
класса Поста, то исключаются функции самобазисы $\downarrow$ и $\phi_{14}$. Функции
$\phi_3$ и $\phi_5$ принадлежат всем классам Поста, поэтому с ними невозможно
образовать минимальную систему, какой является базис. Сами комбинации ищутся перебором
\begin{enumerate}
  \item $\mathcal{B} = \{\land,\,+,\,\iff\}$
  \item $\mathcal{B} = \{\land,\,+,\,\phi_{15}\}$
  \item $\mathcal{B} = \{\land,\,\phi_0,\,\iff\}$
  \item $\mathcal{B} = \{\lor,\,+,\,\iff\}$
  \item $\mathcal{B} = \{\lor,\,+,\,\phi_{15}\}$
  \item $\mathcal{B} = \{\lor,\,\phi_0,\,\iff\}$
\end{enumerate}

\subsection{Выразите все булевы функции двух переменных в базисе $\phi_8$}

\end{document}